O desenvolvimento do HormuzNet permitiu a consolidação de conceitos fundamentais de sistemas distribuídos em uma aplicação prática e crítica. A transição de uma arquitetura centralizada para uma malha descentralizada provou ser eficaz para aumentar a resiliência do monitoramento marítimo no Estreito de Ormuz.

As decisões de design, como o uso de UDP Multicast para aquisição redundante e Relógios de Lamport para consistência de estado, foram validadas pelos testes de estresse e injeção de falhas. O sistema demonstrou capacidade de auto-recuperação e manutenção da integridade operacional mesmo diante da perda de componentes vitais, como \textit{brokers} e drones.

Trabalhos futuros poderiam explorar a integração de algoritmos de eleição de líder para coordenação de tarefas mais complexas e o uso de criptografia na malha P2P para garantir a segurança da comunicação tática. Em suma, o HormuzNet estabelece uma base sólida para sistemas de monitoramento autônomos em ambientes de alta criticidade.
